%==============================================================================
% Appendix C: KK vs Matsubara coexistence tables
%==============================================================================
\section{KK vs.\ Matsubara Coexistence Tables}
\label{app:C}

This appendix tabulates the two readings of the same Euclidean circle.  The
purpose is to make it visible at a glance that the KK reading and the
Matsubara-style thermal reading coexist as secondary readings of the same
Euclidean data.

\subsection{Basic Coexistence Table}
\label{app:C:basic}

\begin{table}[H]
\centering\small
\begingroup
\renewcommand{\arraystretch}{2.0}
\begin{tabularx}{\linewidth}{YYYcY}
\toprule
\shortstack{Same Euclidean\\quantity} & KK reading & \shortstack{Matsubara\\reading} & coexist? & Caution \\
\midrule
$\Sone$ length $L$ & compactification scale & inverse-temperature\newline circle with $\beta = L$ & yes & same circle, two readings \\
Maxwell baseline & reduced kinetic baseline & thermal stiffness scale & yes & shared coefficient\newline does not imply identity \\
CS-active structure & reduced active\newline channel count & thermal parity-odd\newline channel count & yes & interpretation differs \\
$\SCZ$ & reduced inflow action & interface free-energy\newline descriptor & yes & not a real-time statement \\
branch label & reduced branch language & thermal branch language & yes & not a full\newline finite-temperature derivation \\
\bottomrule
\end{tabularx}
\endgroup
\caption{KK vs.\ Matsubara coexistence table.}
\label{tab:appC_basic}
\end{table}

The point of this table is to distinguish ``shared coefficient'' from
``identical theory''.  The fact that the same Euclidean quantity admits two
descriptor languages does not imply complete agreement of the two
theoretical readings.

\subsection{Representative Quantities and the Two Readings}
\label{app:C:rep}

\begin{table}[H]
\centering\small
\begingroup
\renewcommand{\arraystretch}{2.0}
\begin{tabularx}{\linewidth}{YYYY}
\toprule
Native quantity & \shortstack{KK /\\reduced reading} & \shortstack{Matsubara /\\thermal reading} & Remark \\
\midrule
$\KMxw$ & reduced kinetic baseline & thermal stiffness scale & same coefficient,\newline different interpretation \\
CS direction\newline count & reduced parity-odd\newline channel count & thermal parity-odd\newline channel count & counting survives,\newline meaning differs \\
$\etaAPS$ & lower-dimensional\newline spectral asymmetry & thermal spectral descriptor & not a transport coefficient\newline by itself \\
$\Ntop$ & reduced interface CS jump & thermal torsional-sector label & remains frame-bundle-norm.\ torsional charge \\
$\SCZ$ & reduced inflow functional & thermal interface\newline free-energy descriptor & Euclidean, not real-time \\
\bottomrule
\end{tabularx}
\endgroup
\caption{Representative quantities and the two readings.}
\label{tab:appC_rep}
\end{table}

\subsection{Interpretive Distinction and Scope}
\label{app:C:scope}

The KK reading provides a lower-dimensional descriptor language and the
Matsubara reading provides a thermal descriptor language.  The two are based
on the same Euclidean circle, but neither claims theoretical identity with
the other.

The Matsubara reading preserves the original product structure
$M^3\times \Sone$, so within the scope of this paper it furnishes a
$P=0$-protected thermal sector.  The thermal language here therefore does
not include the derivation of full finite-temperature field theory, the
complete analysis of the Matsubara sum, or real-time transport.
